\documentclass[12pt,a4paper]{article}
\usepackage[utf8]{inputenc}
\usepackage{cmap}
\usepackage[T2A]{fontenc}
\usepackage[russian]{babel}
\usepackage{geometry}
\usepackage{booktabs}
\geometry{top=2cm, left=2cm, right=2cm, bottom=2cm}

\title{Лабораторная работа №5. Теория и вопросы защиты: дерево решений для классификации грибов}
\author{Ануфриев А.С. 465029}
\date{08.11.2026}

\begin{document}
\maketitle

\section{Деревья решений}
Дерево решений --- иерархическая структура, представляющая правила классификации в виде последовательности проверок. Каждый внутренний узел содержит проверку по какому-то признаку, ребра соответствуют возможным значениям признака, а листья содержат метки классов. Каждый объект попадает в единственный лист, что даёт покрывающий набор конъюнкций. Классы должны быть дискретными, каждый пример относится к одному классу однозначно.

\section{Алгоритм C4.5}
Алгоритм C4.5 строит небинарное дерево (количество потомков узла не ограничено) и решает задачи классификации. Работа происходит итеративно сверху вниз:
\begin{enumerate}
\item изначально все примеры ассоциированы с корнем;
\item выбирается признак для проверки, множество разбивается на подмножества по значениям признака, для каждого создаётся потомок;
\item процесс рекурсивно повторяется для каждого подмножества до остановки.
\end{enumerate}
Количество классов должно быть значительно меньше количества записей.

\subsection{Критерий разбиения}
Пусть $T$ --- множество примеров, $freq(C_j,T)$ --- число примеров класса $C_j$. Энтропия (средняя информация для определения класса):
\[
Info(T) = -\sum_{j=1}^{k} \frac{freq(C_j,T)}{|T|} \log_2 \frac{freq(C_j,T)}{|T|}.
\]
После разбиения по проверке $X$ со значениями $A_1,\dots,A_n$, дающем подмножества $T_i$:
\[
Info_X(T) = \sum_{i=1}^{n} \frac{|T_i|}{|T|} Info(T_i), \qquad
splitinfo(X) = -\sum_{i=1}^{n} \frac{|T_i|}{|T|} \log_2 \frac{|T_i|}{|T|}.
\]
Нормированный прирост информации:
\[
GainRatio(X) = \frac{Info(T) - Info_X(T)}{splitinfo(X)}.
\]
На каждом шаге выбирается признак с максимальным $GainRatio$.

\section{Реализация без сторонних библиотек}
Дерево реализуется через узел со ссылками на потомка по каждому значению признака (не список-списков). В узле хранятся индекс признака-разделителя и словарь ``значение признака $\rightarrow$ дочерний узел''. Лист хранит класс большинства и долю объектов положительного класса, которая используется как вероятность предсказания. Остановка рекурсии: все примеры в узле одного класса, отсутствие прироста информации, исчерпание признаков или достижение максимальной глубины.

\section{Метрики качества бинарной классификации}
Пусть объект относится к положительному классу (ядовитый гриб). Матрица ошибок: TP (истинноположительные), FP (ложноположительные), TN (истинноотрицательные), FN (ложноотрицательные). Тогда:
\[
Accuracy = \frac{TP+TN}{TP+FP+TN+FN}, \qquad
Precision = \frac{TP}{TP+FP}, \qquad
Recall = \frac{TP}{TP+FN}.
\]
Все метрики вычисляются собственным кодом без библиотек.

\section{Кривые AUC-ROC и AUC-PR}
При изменении порога принятия решения строятся пары $(FPR, TPR)$, где $TPR = TP/(TP+FN)$ и $FPR = FP/(FP+TN)$ --- это ROC-кривая; площадь под ней --- AUC-ROC. Пара $(Recall, Precision)$ даёт PR-кривую, её площадь --- AUC-PR. Площади считаются методом трапеций по отсортированным по убыванию вероятностей точкам. Обе кривые и значения AUC вычисляются без сторонних библиотек.

\section{Вероятные вопросы защиты с ответами}
\begin{enumerate}
\item \emph{Что представляет собой дерево решений и какую задачу оно решает?}
Ответ: иерархическая структура правил классификации; узлы --- проверки признаков, ребра --- значения признаков, листья --- классы; решает задачи классификации с дискретными классами.

\item \emph{Чем алгоритм C4.5 отличается от бинарных деревьев?}
Ответ: количество потомков узла не ограничено, каждое значение признака порождает отдельного потомка.

\item \emph{По какому критерию выбирается признак для разбиения узла?}
Ответ: по максимальному нормированному приросту информации $GainRatio(X)$.

\item \emph{Как считается энтропия множества примеров?}
Ответ: $Info(T) = -\sum_j p_j \log_2 p_j$, где $p_j$ --- доля примеров класса $C_j$.

\item \emph{Что делает множитель splitinfo в gain ratio?}
Ответ: корректирует прирост информации, снижая преимущество признаков с большим числом уникальных значений.

\item \emph{Почему в данных были исключены niektóre записи?}
Ответ: в признаке ``stalk-root'' встречаются пропуски ``?'', такие объекты удалены перед обучением.

\item \emph{Сколько признаков использовалось и как они выбирались?}
Ответ: $\sqrt{22} \approx 4$ признака, отобранных случайным образом с фиксированным зерном генератора.

\item \emph{Как вычислялись Accuracy, precision и recall?}
Ответ: по формулам из матрицы ошибок, собственным кодом без библиотек.

\item \emph{Что такое AUC-ROC и AUC-PR и как они вычислялись?}
Ответ: площади под ROC- и PR-кривыми, построенными по парам (FPR, TPR) и (Recall, Precision) при разных порогах; площадь считается методом трапеций, без библиотек.

\item \emph{Какое качество достигнуто на тестовой выборке?}
Ответ: Accuracy 0.8547, precision 0.8692, recall 0.7307, AUC-ROC 0.9406, AUC-PR 0.9223.
\end{enumerate}

\end{document}
